\subsection{Limitations and threats to validity}
\label{sec:limitations}

\noindent\textbf{Internal validity.} The main tables choose the 1\%-FPR threshold on the evaluated benign set, which favors every detector. Section~\ref{sec:robust} replaces it with thresholds calibrated on held-out suites; the ranking at the top holds, but individual numbers drop, and the held-out numbers are the ones a deployment should expect. The benign sets are small at the scale of a 1\% budget: AD-Matched has \NADBenign benign outputs, so its 1\%-FPR threshold is set by about six of them, and the Llama judge's 95\% interval on AD-Matched spans \ADJLTprCI\% (Table~\ref{tab:main}). Larger benign sets, which require more agent tasks rather than more attacks, would tighten these estimates.

\noindent\textbf{Construct validity.} Replay executes the ground-truth calls of each task, whereas a real agent may call other tools, call them with other arguments, or fail. Section~\ref{sec:robust} shows that replay reproduces what a real agent receives for the same calls, up to write timestamps, and that output-level false-positive rates carry over to a real agent's outputs; this check covers one agent model and benign runs only. On $\tau$-bench we chose the injection point (product names) and the attack set (InjecAgent) ourselves, and another placement could favor different detectors. Whether an injected output would have misled an agent is not measured, since a detector must flag the text before any agent reads it.

\noindent\textbf{External validity.} Both agent benchmarks are simulations that return synthetic YAML or JSON records. Real tool outputs, such as HTML pages, API responses, and long documents, are messier, and format alone moves false-positive rates, so absolute numbers will differ in deployment; we expect the relative findings that held on both benchmarks to carry over. All injections come from benchmark templates, and an adaptive attacker who knows the detector or judge would lower every detection rate~\cite{attackermovessecond,zhan2025adaptive,hackett2025}, so our detection numbers are upper bounds, whereas the false-positive results do not depend on the attacker. We did not run adaptive or paraphrasing attacks of our own: such an evaluation is a study of attacks that needs its own threat model and attack budget, and combining it with our replay methodology is future work. Commercial API detectors are not included, and only PIGuard and Horizon-Labs publish enough to audit their training data. Our audits find verbatim copies only; for PIGuard and BIPIA, verbatim overlap alone already explains the score.

\noindent\textbf{Statistical conclusion validity.} With \NDetectors detectors, rank correlations have little power, and a non-significant correlation does not show that two rankings are unrelated. The claim that rankings do not transfer from BIPIA to AgentDojo therefore rests on concrete reversals (PIGuard: \BIPGTpr on BIPIA, \ADPGTpr on AgentDojo), not on the near-zero $\tau$, and between the agent benchmarks we claim no more than partial agreement ($\tau=\TauADTB$, $p=\TauADTBP$). The false-positive correlation is the one rank result that is significant under both the standard and the permutation test.

\noindent\textbf{Judges as a diagnostic.} The judges use a single prompt about actions and were not tested against adaptive attacks; a prompt that also asks about changes to the answer might cover more attack goals, at an unmeasured cost in false alarms.
